% =====================================================================
%  Détection de dessins doublons non triviaux — bibliographie, mesures
%  et spécification d'une méthode (media_restorer)
%
%  Compilation :  pdflatex rapport-doublons-dessins.tex   (× 3)
% =====================================================================
\documentclass[11pt,a4paper]{article}

\usepackage[utf8]{inputenc}
\usepackage[T1]{fontenc}
\usepackage{lmodern}
\usepackage[french]{babel}
\usepackage[margin=2.4cm]{geometry}
\usepackage{microtype}
\usepackage{parskip}
\usepackage{booktabs}
\usepackage{tabularx}
\usepackage{enumitem}
\usepackage{xcolor}
\usepackage{amsmath}
\usepackage{fancyhdr}
\usepackage{titlesec}
\usepackage[hidelinks]{hyperref}
\usepackage{xurl}             % césure des URL n'importe où

\definecolor{accent}{HTML}{1F4E79}
\definecolor{encadre}{HTML}{F2F5F8}
\definecolor{alerte}{HTML}{A33A2E}
\definecolor{vert}{HTML}{2E6B3E}
\hypersetup{
  colorlinks=true, linkcolor=accent, urlcolor=accent, citecolor=accent,
  pdftitle={Détection de dessins doublons non triviaux},
  pdfsubject={media\_restorer},
  pdfkeywords={doublons, copy detection, homographie, MAGSAC, invariance},
}

\titleformat{\section}{\Large\bfseries\color{accent}}{\thesection}{0.7em}{}
\titleformat{\subsection}{\large\bfseries}{\thesubsection}{0.6em}{}

\pagestyle{fancy}
\fancyhf{}
\fancyhead[L]{\small\color{accent}media\_restorer — doublons}
\fancyhead[R]{\small\thepage}
\renewcommand{\headrulewidth}{0.4pt}

\newcommand{\aretenir}[1]{%
  \par\medskip
  \noindent\colorbox{encadre}{%
    \begin{minipage}{\dimexpr\linewidth-2\fboxsep}
      \vspace{2pt}#1\vspace{2pt}
    \end{minipage}}%
  \par\medskip
}
\newcommand{\alerte}[1]{\textcolor{alerte}{\textbf{#1}}}
\newcommand{\bon}[1]{\textcolor{vert}{\textbf{#1}}}

\newcolumntype{Y}{>{\raggedright\arraybackslash}X}
% Colonne FIXE en drapeau. Une colonne « l » ne se coupe jamais : y mettre
% du texte long lui fait absorber toute la largeur du tableau, et les
% colonnes X restantes débordent dans la marge.
\newcolumntype{F}[1]{>{\raggedright\arraybackslash}p{#1}}

% =====================================================================
\begin{document}

\begin{center}
  {\LARGE\bfseries\color{accent} Détection de dessins doublons\\[2pt]
   non triviaux\par}
  \vspace{6pt}
  {\large Bibliographie, mesures et spécification d'une méthode\par}
  \vspace{4pt}
  {\small Projet \texttt{media\_restorer} \textbullet{} corpus
   \texttt{Dessins/Cabrol} : 8\,693 images\par}
  {\small Toutes les mesures de ce rapport ont été faites sur ce corpus\par}
  {\small \today\par}
\end{center}

\vspace{4pt}\hrule\vspace{10pt}
\tableofcontents
\vspace{10pt}\hrule

% =====================================================================
\section{Le mur combinatoire}
\label{sec:mur}

Comparer toutes les paires du corpus, c'est \textbf{37\,779\,778 comparaisons}.
Les coûts ci-dessous ne sont pas des estimations : ils ont été mesurés sur des
dessins du corpus, en travaillant à 1\,000\,px de côté (voir
§\ref{sec:mesures}).

\begin{table}[h]
\centering\small
\begin{tabular}{lrrr}
\toprule
\textbf{Détecteur} & \textbf{Coût mesuré} & \textbf{Toutes les paires}
  & \textbf{Top-20 candidats} \\
\midrule
ORB   & \phantom{0}18 ms & \alerte{7,9 jours}  & \bon{0,43 h} \\
AKAZE & \phantom{0}71 ms & \alerte{31,0 jours} & \bon{1,71 h} \\
SIFT  & 148 ms           & \alerte{64,7 jours} & \bon{3,57 h} \\
\bottomrule
\end{tabular}
\caption{Vérification géométrique exhaustive contre vérification ciblée sur
20 candidats par image (86\,930 paires). Rapport : \textbf{434×}.}
\end{table}

\aretenir{\textbf{Une architecture à deux étages n'est pas un raffinement, c'est
la condition d'existence du projet.} Il faut d'abord \emph{proposer} des
candidats à bas coût, puis ne \emph{vérifier} géométriquement que ceux-là.}

\subsection{Un résultat contre-intuitif : aucun index approché n'est nécessaire}

L'étage « candidats » se ramène à comparer 8\,693 descripteurs entre eux. Avec
des vecteurs de 256 dimensions, la matrice de similarité complète représente
\textbf{19,3 G opérations multiplication-accumulation} et \textbf{302\,Mo} —
soit \textbf{quelques secondes} avec BLAS sur les 16 c\oe urs de la machine, et
une empreinte négligeable devant les 92\,Go de RAM disponibles.

\alerte{Aucune bibliothèque d'index approché n'est donc requise.}
\texttt{faiss}, \texttt{hnswlib} et \texttt{annoy} sont absents de
l'installation, et le resteraient utilement : ils ne deviennent nécessaires
qu'au-delà de $10^5$ images environ. À 8\,693, la force brute est à la fois plus
simple et \emph{exacte}, là où un index approché introduirait des faux négatifs
sans contrepartie.

% =====================================================================
\section{Trois régimes, et non un problème unique}
\label{sec:regimes}

Les quatre situations à détecter ne relèvent pas de la même mécanique. La
question qui les sépare est simple : \emph{existe-t-il une transformation
géométrique qui envoie une image sur l'autre ?}

\begin{table}[h]
\centering\small
\begin{tabularx}{\textwidth}{llY}
\toprule
\textbf{Régime} & \textbf{Situation} & \textbf{Une homographie relie-t-elle les deux images ?} \\
\midrule
\textbf{R1} copie géométrique & même original, deux numérisations
  & \bon{oui}, globale — les deux images se recouvrent entièrement \\
\addlinespace
\textbf{R2} inclusion partielle & même dessin sur un autre support ;
  détail, recadrage
  & \bon{oui}, mais vers une \textbf{sous-région} : la transformation existe,
    le recouvrement est partiel \\
\addlinespace
\textbf{R3} variante redessinée & le même sujet redessiné par l'artiste
  & \alerte{non} — aucune transformation ne relie les traits \\
\bottomrule
\end{tabularx}
\end{table}

\aretenir{\textbf{R1 et R2 se traitent par la même machinerie} — l'inclusion
partielle n'est pas un cas particulier à ajouter, elle sort naturellement de
l'estimation (§\ref{sec:dico}). \alerte{R3 en est fondamentalement distinct} :
aucune méthode géométrique ne peut l'atteindre, et les méthodes qui le peuvent
sont nettement moins fiables. Présenter une « solution unique » aux quatre cas
serait malhonnête.}

% =====================================================================
\section{Les familles de méthodes}
\label{sec:familles}

\begin{table}[h]
\centering\small
\begin{tabularx}{\textwidth}{F{2.5cm}YYF{2.1cm}}
\toprule
\textbf{Famille} & \textbf{Invariances} & \textbf{Rôle} & \textbf{Coût} \\
\midrule
\textbf{A} — hachage perceptuel
  (pHash, dHash) & contraste, translation faible
  & \alerte{pré-filtre quasi-exact seulement} : ni rotation, ni recadrage
  & négligeable \\
\addlinespace
\textbf{B} — Fourier-Mellin log-polaire ; moments de Zernike, de Hu
  & rotation, homothétie, translation
  & candidats R1 ; \textbf{estime déjà} angle et échelle
  & faible \\
\addlinespace
\textbf{C} — traits locaux + vérification géométrique
  (SIFT, AKAZE, ORB + MAGSAC) & affine, partielle, photométrique
  & \bon{la réponse pour R1 et R2} & 18--148 ms/paire \\
\addlinespace
\textbf{D} — descripteurs appris (SSCD, DINOv2, CLIP)
  & apprises, larges & candidats ; \textbf{seul espoir pour R3} & moyen \\
\addlinespace
\textbf{E} — spécifique dessin (squelettisation, descripteurs de traits)
  & épaisseur de trait & normalisation \emph{en amont} & faible \\
\bottomrule
\end{tabularx}
\end{table}

\subsection{A — hachage perceptuel : à ne pas confondre avec une solution}

Le hachage perceptuel réduit une image à quelques dizaines de bits et compare
par distance de Hamming. C'est la première idée qui vient, et elle ne répond
pas à la question posée : \alerte{pHash n'est invariant ni à la rotation ni au
recadrage.} Il reste utile comme filtre de premier passage pour les doublons
\emph{quasi} exacts (recompression, changement de format), pas au-delà.

\subsection{B — Fourier-Mellin : la rotation devient une translation}

La méthode de \textsc{Reddy} et \textsc{Chatterji} exploite une propriété
élégante : dans le domaine de Fourier, une translation ne change que la phase,
donc le \emph{module} du spectre est invariant par translation. En passant ce
module en coordonnées log-polaires, une rotation et une homothétie deviennent
deux \emph{translations} — que l'on retrouve par corrélation de phase.

Intérêt pour l'étage candidats : le descripteur est global (donc rapide), et il
fournit \textbf{déjà} une estimation de l'angle et du facteur d'échelle.
\texttt{cv2.warpPolar} est disponible dans l'installation : rien à ajouter.
Limite : étant \emph{global}, il tolère mal l'inclusion partielle — il sert donc
R1, pas R2.

\subsection{C — traits locaux et vérification géométrique}

Le schéma classique : détecter des points d'intérêt, les décrire, apparier par
plus proche voisin avec le test de ratio de \textsc{Lowe}, puis estimer une
homographie robuste qui rejette les appariements aberrants. Le nombre
d'appariements survivants — les \emph{inliers} — est la mesure de similarité.

\bon{C'est cette famille qui répond à R1 et R2}, et l'inclusion partielle y est
native : rien n'oblige les points appariés à couvrir toute l'image.

Deux choix modernes importent. Le premier est l'estimateur robuste :
\textbf{MAGSAC++} (\texttt{cv2.\allowbreak USAC\_\allowbreak MAGSAC}, présent dans OpenCV 4.13) évite le
réglage délicat du seuil d'inlier de RANSAC. Le second est le détecteur, que la
§\ref{sec:mesures} tranche par la mesure plutôt que par réputation.

\subsection{D — descripteurs appris}

L'état de l'art de la détection de copie est aujourd'hui \textbf{SSCD}
(CVPR 2022), entraîné par apprentissage auto-supervisé contrastif et évalué sur
le jeu \textbf{DISC21} du \emph{Image Similarity Challenge} de Meta — c'est-à-dire
exactement notre problème : retrouver une image transformée parmi un million de
distracteurs. SSCD y dépasse les descripteurs auto-supervisés généralistes de
très loin (jusqu'à +48 points absolus sur SimCLR).

Pour les variantes redessinées (R3), la littérature récente sur la recherche
d'\oe uvres d'art situe \textbf{DINOv2} devant CLIP en similarité fine d'images
— CLIP restant préférable si l'on veut interroger par texte. \texttt{transformers}
étant déjà une dépendance du projet, CLIP est utilisable \emph{ce soir} ;
DINOv2 et SSCD demanderaient un téléchargement de poids, mais \textbf{aucun
entraînement} — ce qui les met à portée de la machine malgré l'impossibilité
d'entraîner sur le GPU 780M.

\subsection{E — l'épaisseur de trait : un axe qu'aucune autre famille ne couvre}
\label{sec:trait}

Aucune des familles A à D n'est conçue pour ignorer l'épaisseur du trait. Or
c'est justement l'axe sur lequel les mesures (§\ref{sec:mesures}) montrent la
plus grande fragilité. Deux remèdes :

\begin{itemize}[leftmargin=*,itemsep=3pt]
  \item \textbf{Squelettiser avant de décrire} — ramener tout trait à une
  largeur d'un pixel rend la description indépendante de l'épaisseur.
  \alerte{\texttt{cv2.ximgproc.thinning} est absent} de l'installation, mais
  \textbf{le projet sait déjà extraire un squelette} :
  \texttt{engines/vectorise/topology.py} le fait par homologie persistante
  (GUDHI), sans \texttt{scikit-image}.
  \item \textbf{Normaliser morphologiquement} — égaliser les épaisseurs
  médianes des deux images (mesurées par transformée de distance) avant
  comparaison. Moins radical, beaucoup moins coûteux.
\end{itemize}

% =====================================================================
\section{Mesures sur le corpus}
\label{sec:mesures}

Douze dessins tirés au hasard dans \texttt{Cabrol}, ramenés à 1\,000\,px de
côté. Chacun est comparé à sa propre version transformée : la vérité terrain est
donc \emph{parfaite}, on sait exactement quelle transformation a été appliquée.
La grandeur relevée est la \textbf{médiane du ratio d'inliers} — la proportion
d'appariements survivant à la vérification géométrique.

\begin{table}[h]
\centering\small
\begin{tabular}{lccc}
\toprule
\textbf{Transformation} & \textbf{SIFT} & \textbf{AKAZE} & \textbf{ORB} \\
\midrule
identité                  & 1,00 & 1,00 & 1,00 \\
rotation 5°               & 0,96 & 0,98 & 0,95 \\
rotation 15°              & 0,97 & 0,98 & 0,95 \\
rotation 45°              & 0,97 & 0,96 & 0,95 \\
rotation 90°              & 1,00 & 0,99 & 0,94 \\
homothétie ×0,5           & 0,88 & 0,91 & \bon{0,98} \\
homothétie ×2             & 0,97 & 0,97 & 0,85 \\
contraste +40\,\%         & 0,99 & 0,99 & 1,00 \\
contraste $-$40\,\%         & 0,99 & 0,94\,\textsuperscript{\dag} & 0,99 \\
\alerte{trait +2 px}      & 0,85 & 0,95 & 0,93 \\
\alerte{trait $-$2 px}      & \alerte{0,76}\,\textsuperscript{\dag} & 0,89\,\textsuperscript{\dag} & 0,90 \\
recadrage à 25\,\%        & 0,96 & 0,95 & 0,95 \\
rot. 15° + ×0,7 + contraste & 0,94 & 0,95 & 0,97 \\
JPEG qualité 40           & 0,97 & 0,99 & 0,99 \\
\midrule
\textbf{coût médian}      & 148 ms & 71 ms & \bon{18 ms} \\
\bottomrule
\end{tabular}
\caption{Médiane du ratio d'inliers sur 12 dessins.
\textsuperscript{\dag} au moins un dessin est descendu à un ratio quasi nul
(SIFT 0,25 ; AKAZE 0,00) — voir ci-dessous.}
\end{table}

\subsection{Trois enseignements}

\aretenir{\textbf{(1) ORB suffit, et il est huit fois plus rapide que SIFT.}
Sur ces dessins, ORB tient la rotation (0,94--0,95 quel que soit l'angle), le
contraste, le recadrage et la compression aussi bien que SIFT, pour 18\,ms
contre 148. C'est \emph{contraire} à la réputation d'ORB, établie sur des
photographies : le trait au noir sur fond clair produit des coins très
contrastés, terrain idéal pour un détecteur FAST.}

\textbf{(2) L'épaisseur de trait est le point faible, pour les trois.} C'est le
seul axe où la médiane descend sous 0,90, et surtout le seul où l'on observe des
\alerte{effondrements individuels} : un dessin est tombé à 0,25 avec SIFT, un
autre à 0,00 avec AKAZE — c'est-à-dire une non-détection complète. L'intuition
de départ était donc juste, et elle justifie à elle seule la famille E
(§\ref{sec:trait}).

\textbf{(3) ORB et l'homothétie : attention au sens.} ORB est excellent en
réduction (0,98 à ×0,5) et faiblit en agrandissement (0,85 à ×2). En pratique on
peut toujours \emph{réduire} la plus grande des deux images plutôt que d'agrandir
la plus petite — ce qui place systématiquement ORB dans son bon régime, et coûte
moins cher au passage.

% =====================================================================
\section{Le dictionnaire de mérite}
\label{sec:dico}

\aretenir{\textbf{Idée directrice : la transformation estimée \emph{est}
l'explication.} On ne construit pas un score explicable après coup — on lit
l'explication dans l'homographie que la vérification a déjà calculée.}

\subsection{Décomposer l'homographie}

Soit $H$ l'homographie estimée et $A$ sa partie linéaire $2\times2$. Sa
décomposition en valeurs singulières $A = U\Sigma V^{\mathsf T}$ donne
immédiatement :

\begin{itemize}[leftmargin=*,itemsep=2pt]
  \item la \textbf{rotation} : l'angle de la matrice orthogonale $UV^{\mathsf T}$ ;
  \item l'\textbf{échelle} : $\sqrt{\sigma_1\sigma_2}$ ;
  \item le \textbf{cisaillement} : l'anisotropie $\sigma_1/\sigma_2$, égale à 1
    pour une similitude pure ;
  \item la \textbf{translation} : la dernière colonne de $H$.
\end{itemize}

\textbf{Vérifié expérimentalement} — on applique une transformation connue, puis
on regarde ce que la décomposition en retrouve :

\begin{table}[h]
\centering\small
\begin{tabular}{lrrr}
\toprule
\textbf{Appliqué} & \textbf{Rotation estimée} & \textbf{Échelle estimée}
  & \textbf{Anisotropie} \\
\midrule
identité                  & \phantom{$-$}0,00° & 1,000 & 1,0000 \\
rotation 15°              & 15,00°           & 1,000 & 1,0002 \\
rotation 45°              & 44,98°           & 1,000 & 1,0008 \\
échelle 0,60              & \phantom{$-$}0,01° & 0,600 & 1,0011 \\
rot. 30° + échelle 0,75   & 30,01°           & 0,750 & 1,0013 \\
rot. $-$20° + échelle 1,50  & 20,01°           & 1,500 & 1,0006 \\
\bottomrule
\end{tabular}
\caption{Médianes sur 8 dessins. Erreur de rotation $<0{,}02°$, erreur d'échelle
$<0{,}001$. Le signe de l'angle dépend de la convention d'axes (l'axe $y$ est
orienté vers le bas en coordonnées image) : c'est une convention à fixer, pas
une erreur.}
\end{table}

\subsection{La couverture asymétrique détecte l'inclusion partielle}

C'est le mécanisme qui répond au cas « une partie d'une image contient une autre
image ». On projette les quatre coins de $A$ dans le repère de $B$ par $H$, et
l'on mesure l'aire de l'intersection rapportée à celle de $B$ ; puis la même
chose dans l'autre sens avec $H^{-1}$. \textbf{Les deux valeurs diffèrent, et
c'est précisément l'information utile.}

\begin{table}[h]
\centering\small
\begin{tabular}{rrrr}
\toprule
\textbf{Part de la page occupée} & \textbf{Inliers} &
\textbf{Couverture dessin$\to$page} & \textbf{Couverture page$\to$dessin} \\
\midrule
50\,\% & 1\,362 & 0,50 & 1,00 \\
25\,\% & \phantom{0}750 & 0,25 & 1,00 \\
10\,\% & \phantom{0}248 & 0,10 & 1,00 \\
\phantom{0}4\,\% & \phantom{00}12 & 0,04 & 1,00 \\
\bottomrule
\end{tabular}
\caption{Un dessin collé dans une page plus grande. La couverture vaut
\textbf{exactement} la fraction d'aire occupée, et la détection survit jusqu'à
4\,\% de la page (12 inliers).}
\end{table}

\aretenir{\textbf{Deux couvertures égales et proches de 1} $\Rightarrow$ même
dessin (R1). \textbf{Une couverture proche de 1, l'autre faible} $\Rightarrow$
l'un est contenu dans l'autre (R2), et la valeur faible \emph{dit quelle
proportion} il occupe. Un scalaire unique de similarité ne pourrait pas
distinguer ces deux situations.}

\subsection{Schéma proposé}

\begin{verbatim}
{
  "regime": "geometrique" | "partielle" | "semantique",
  "merite": 0.0-1.0,              # score agrégé, comparable entre régimes

  # --- peuplés en R1 / R2 seulement ---------------------------------
  "n_inliers": int,               "ratio_inliers": float,
  "rotation_deg": float,          "echelle": float,
  "anisotropie": float,           "translation_px": (float, float),
  "couverture_a_dans_b": float,   "couverture_b_dans_a": float,
  "gain_photometrique": float,    "offset_photometrique": float,
  "ratio_epaisseur_trait": float,

  # --- toujours peuplé ----------------------------------------------
  "cosinus_semantique": float,
}
\end{verbatim}

Les champs photométriques s'obtiennent après recalage, en ajustant
$I_B \approx a\,I_B + b$ sur les pixels appariés : $a$ mesure l'écart de
contraste, $b$ celui de luminosité. Le ratio d'épaisseur se mesure de même,
\emph{après} recalage, par transformée de distance sur les deux masques d'encre.

\alerte{Un dictionnaire partiellement vide est une information, pas un défaut} :
en R3, l'absence de champs géométriques \emph{signifie} qu'aucune transformation
n'a été trouvée, et c'est exactement ce que l'utilisateur doit savoir avant de
faire confiance au score.

% =====================================================================
\section{Architecture recommandée}
\label{sec:archi}

Quatre étages, du gratuit au coûteux.

\begin{enumerate}[leftmargin=*,itemsep=5pt]
  \item \textbf{Étage 0 — pré-filtre gratuit.} Réutiliser les mesures
  \textbf{déjà calculées et mises en cache} par
  \texttt{media\_restorer.engines.triage} : orientation, densité d'encre,
  saturations d'encre et de papier. Deux dessins dont la densité d'encre diffère
  d'un facteur trois ne sont pas des doublons. Coût nul, sur des données
  existantes.

  \item \textbf{Étage 1 — candidats.} Descripteur global (Fourier-Mellin, et/ou
  descripteur appris), matrice de similarité en force brute, top-20 par image.
  Quelques secondes (§\ref{sec:mur}).

  \item \textbf{Étage 2 — vérification.} ORB + \texttt{cv2.USAC\_MAGSAC} sur les
  86\,930 paires retenues, produisant le dictionnaire du §\ref{sec:dico}.
  \textbf{26 minutes} au coût mesuré. Repasser en SIFT sur les seuls cas
  ambigus reste abordable (3,6\,h pour tout, donc quelques minutes pour un
  résidu).

  \item \textbf{Étage 3 — repli R3.} Pour les paires que la géométrie rejette,
  cosinus sémantique seul, \textbf{explicitement marqué peu fiable}.
\end{enumerate}

\aretenir{Budget total réaliste : \textbf{de l'ordre de l'heure} pour tout le
corpus, sur le seul CPU, sans rien installer d'autre que d'éventuels poids de
modèle. C'est le rapport 434× entre vérification exhaustive et ciblée
(§\ref{sec:mur}) qui rend le projet possible.}

% =====================================================================
\section{Faisabilité sur la machine}

\begin{table}[h]
\centering\small
\begin{tabularx}{\textwidth}{lY}
\toprule
\textbf{Disponible} & SIFT, ORB, AKAZE, BRISK, KAZE, \texttt{USAC\_MAGSAC},
\texttt{warpPolar} (OpenCV \textbf{4.13}) ; \texttt{transformers} 4.57.6 pour
CLIP ; \texttt{numpy}, \texttt{scipy} \\
\addlinespace
\textbf{Absent} & \texttt{faiss}, \texttt{hnswlib}, \texttt{annoy}
(\bon{inutiles} à cette échelle, §\ref{sec:mur}) ; \texttt{sklearn} ;
\texttt{imagehash} ; \texttt{cv2.ximgproc} (donc pas de \texttt{thinning} —
voir §\ref{sec:trait}) \\
\bottomrule
\end{tabularx}
\end{table}

Les familles A, B, C et E ne demandent \textbf{aucune installation}. La
famille D exige de télécharger des poids, mais \textbf{aucun entraînement} :
l'impossibilité d'entraîner sur le GPU 780M (voir
\texttt{rapport-reperage-automatique-dessins.tex}, §2) n'est donc pas un
obstacle ici.

% =====================================================================
\section{Validation}
\label{sec:validation}

Deux jeux complémentaires, l'un contrôlé, l'autre réel.

\textbf{Synthétique} — les transformations du §\ref{sec:mesures}. Vérité terrain
parfaite, invariance mesurable axe par axe. C'est le seul moyen de savoir
\emph{quelle} différence fait échouer une méthode.

\textbf{Paires réelles} — le corpus contient \textbf{1\,587 noms de fichiers
présents dans au moins deux dossiers différents}. Je craignais une vérité
terrain très bruitée, une même numérotation d'appareil pouvant produire des
homonymes sans rapport. \textbf{La mesure dit le contraire} :

\begin{table}[h]
\centering\small
\begin{tabular}{lr}
\toprule
Paires de même nom testées & 25 \\
Reconnues comme vrais doublons (ratio d'inliers $\geq0{,}5$) & \bon{25} \\
Collisions de nom & 0 \\
Ratio d'inliers observé & 0,87 à 1,00 \\
\bottomrule
\end{tabular}
\end{table}

\alerte{Nuance statistique nécessaire} : 25 paires sur 1\,587, c'est un
échantillon étroit. Zéro échec sur 25 est compatible avec un taux réel de
collisions allant jusqu'à environ 14\,\% (intervalle de confiance à 95\,\%). La
conclusion prudente est donc « les homonymes sont \emph{majoritairement} de
vrais doublons », et non « il n'y a aucune collision ». Un contrôle sur une
centaine de paires resserrerait utilement cette borne.

% =====================================================================
\section{Limites assumées}

\begin{itemize}[leftmargin=*,itemsep=4pt]
  \item \alerte{R3 restera nettement moins fiable} que R1 et R2. Une variante
  redessinée ne se distingue d'un autre dessin du même auteur, dans le même
  style, sur le même sujet, que par une appréciation qu'aucun descripteur ne
  formalise proprement aujourd'hui. Le rapport propose une piste, pas une
  garantie.

  \item \textbf{Aucun seuil unique ne séparera proprement} doublons et
  non-doublons — c'est précisément pourquoi le mérite est gradué et accompagné
  de son dictionnaire. Le seuil dépend de ce que l'on cherche : dédoublonner un
  fonds ou retrouver toutes les réutilisations d'un dessin n'appellent pas la
  même tolérance.

  \item \textbf{Une revue humaine restera nécessaire.} Sur 8\,693 images, même
  une précision de 99\,\% laisse des centaines de paires à trancher. C'est
  l'argument décisif en faveur d'un mérite \emph{explicable} : « rotation de
  15°, échelle 0,7, couverture 0,25 » se vérifie d'un coup d'\oe il, « score
  0,83 » ne se vérifie pas.

  \item Les mesures portent sur \textbf{12 dessins} pour l'invariance et
  \textbf{8} pour la décomposition. Les écarts entre détecteurs y sont nets,
  mais un corpus plus large pourrait déplacer les valeurs de quelques
  centièmes.
\end{itemize}

% =====================================================================
\section*{Sources}
\addcontentsline{toc}{section}{Sources}

\small\raggedright
\begin{itemize}[leftmargin=*,itemsep=3pt]
  \item E.~Pizzi \emph{et al.}, \emph{A Self-Supervised Descriptor for Image
  Copy Detection} (SSCD), CVPR 2022 —
  \url{https://arxiv.org/abs/2202.10261} ;
  code : \url{https://github.com/facebookresearch/sscd-copy-detection}
  \item M.~Douze \emph{et al.}, \emph{Results and findings of the 2021 Image
  Similarity Challenge} (jeu DISC21) —
  \url{https://arxiv.org/pdf/2202.04007}
  \item B.~S.~Reddy, B.~N.~Chatterji, \emph{An FFT-based Technique for
  Translation, Rotation and Scale-Invariant Image Registration} —
  présentation et implémentation :
  \url{https://sthoduka.github.io/imreg_fmt/}
  \item P.~Lindenberger \emph{et al.}, \emph{LightGlue: Local Feature Matching
  at Light Speed}, ICCV 2023 —
  \url{https://github.com/cvg/LightGlue}
  \item M.~Oquab \emph{et al.}, \emph{DINOv2: Learning Robust Visual Features
  without Supervision} —
  \url{https://arxiv.org/html/2304.07193v2}
  \item \emph{A Review on Near-Duplicate Detection of Images using Computer
  Vision Techniques} — \url{https://arxiv.org/pdf/2009.03224}
  \item \emph{Benchmarking Pretrained Vision Embeddings for Near- and Duplicate
  Detection} — \url{https://arxiv.org/pdf/2312.07273}
  \item Littérature « sketch » : \emph{Sketch retrieval via local dense stroke
  features} —
  \url{https://www.sciencedirect.com/science/article/abs/pii/S0262885615001389}
\end{itemize}

\medskip
\textbf{Mesures.} Tous les chiffres des sections \ref{sec:mur},
\ref{sec:mesures}, \ref{sec:dico} et \ref{sec:validation} ont été obtenus sur le
corpus \texttt{Dessins/Cabrol} avec OpenCV 4.13, sur le CPU seul. Aucune valeur
n'est reprise de la littérature sans être signalée comme telle.

\end{document}
